\documentclass[11pt]{article}

\usepackage[margin=1in]{geometry}
\usepackage[T1]{fontenc}
\usepackage{lmodern}
\usepackage{microtype}
\usepackage{amsmath,amssymb,amsthm,mathtools}
\usepackage{mathrsfs}
\usepackage{booktabs}
\usepackage{array}
\usepackage{enumitem}
\usepackage{hyperref}
\usepackage{xcolor}
\usepackage{fancyhdr}
\usepackage{lastpage}
\usepackage{csquotes}

\definecolor{linkblue}{RGB}{32,83,117}
\hypersetup{
  colorlinks=true,
  linkcolor=linkblue,
  citecolor=linkblue,
  urlcolor=linkblue,
  pdftitle={The Hyperfinite Oracles of DATA MIND 3.2},
  pdfauthor={Brian Tenneson}
}

\pagestyle{fancy}
\fancyhf{}
\lhead{The Hyperfinite Oracles of DATA MIND 3.2}
\rhead{Brian Tenneson}
\cfoot{\thepage\ of \pageref{LastPage}}
\setlength{\headheight}{14pt}

\newtheorem{definition}{Definition}[section]
\newtheorem{proposition}[definition]{Proposition}
\newtheorem{remark}[definition]{Remark}
\newtheorem{principle}[definition]{Principle}

\newcommand{\st}{\operatorname{st}}
\newcommand{\Omn}{\operatorname{Omn}}
\newcommand{\Cn}{\operatorname{Cn}}
\newcommand{\AMLD}{\operatorname{AMLD}}
\newcommand{\Role}{\operatorname{Role}}
\newcommand{\cost}{\operatorname{cost}}
\newcommand{\BANK}{\mathrm{BANK}}
\newcommand{\FBANK}{\mathrm{FUTUREBANK}}
\newcommand{\V}{\mathrm{V}}
\newcommand{\Pcal}{\mathcal{P}}
\newcommand{\Rcal}{\mathcal{R}}
\newcommand{\Ical}{\mathcal{I}}
\newcommand{\Ccal}{\mathcal{C}}
\newcommand{\Ocal}{\mathcal{O}}
\newcommand{\Kcal}{\mathcal{K}}
\newcommand{\Fcal}{\mathcal{F}}

\title{\textbf{The Hyperfinite Oracles of DATA MIND 3.2}\\[0.35em]
\large Logical Omniscience, Nonstandard Horizons, and Verifier-Gated Epistemic Guidance}
\author{Brian Tenneson}
\date{September 2026}

\begin{document}
\maketitle

\begin{center}
\begin{minipage}{0.84\textwidth}
\centering\itshape
``The oracle may direct attention; it may not manufacture certified knowledge.''
\end{minipage}
\end{center}

\begin{abstract}
Standard epistemic logic idealizes knowledge as deductively closed: if an agent knows premises and knows that they imply a conclusion, the agent knows the conclusion. Combined with necessitation, this produces the familiar phenomenon of logical omniscience. DATA MIND, by contrast, is a settlement architecture whose agents must spend finite resources to discover verifier-acceptable certificates. This note develops a nonstandard bridge between the ideal semantic picture and the computational one. A resource-indexed modality $K_i^{\le r}\varphi$ is extended, conceptually, to an unlimited hypernatural horizon $H$. The resulting internal hyperfinite family supports a hyperfinite conjunction of epistemic obligations and a hypernatural closure-cost functional $\Ocal_i(H)$. Standard traces and, where a bounded hyperreal normalization exists, standard-part shadows provide ways to compare the hyperfinite construction with ordinary mathematics.

The paper then introduces a hyperfinite settlement oracle $\mathscr O_H$ with values in the AMLD settlement roles $\{P,R,I,C,U\}$. The oracle is not a verifier and does not possess authority to place propositions in the verified BANK. Instead it may be used as hidden ground truth or, in explicitly counterfactual experiments, as a source of role information that changes awareness and resource routing. DATA MIND 3.2 implements only finite-horizon proxies of these ideas; it does not claim to instantiate an unlimited hypernatural on a conventional computer. The resulting architecture gives a controlled way to measure how much of settlement difficulty comes from recognizing the epistemic role of a problem, as opposed to executing the relevant proof, refutation, independence, or contradiction search once that role is known.
\end{abstract}

\tableofcontents
\newpage

\section{From logical omniscience to resource-indexed knowledge}

Let $K_i\varphi$ mean that epistemic agent $i$ knows $\varphi$. In a normal modal epistemic logic, the distribution axiom is
\begin{equation}
K_i(\varphi\rightarrow\psi)\rightarrow(K_i\varphi\rightarrow K_i\psi),
\label{eq:K}
\end{equation}
and necessitation has the rule form
\begin{equation}
\frac{\vdash\varphi}{\vdash K_i\varphi}.
\label{eq:nec}
\end{equation}
If $\Gamma=\{\varphi_1,\ldots,\varphi_n\}$ and $\Gamma\vdash\psi$, then a corresponding theorem of the object logic can be written in nested implication form. Necessitation followed by repeated applications of \eqref{eq:K} yields
\begin{equation}
\bigl(K_i\varphi_1\wedge\cdots\wedge K_i\varphi_n\bigr)\Longrightarrow K_i\psi.
\end{equation}
Thus an ideal epistemic agent is closed under logical consequence. If it knows the members of $\Gamma$, it knows every member of $\Cn(\Gamma)$.

That abstraction suppresses search cost. DATA MIND is built precisely in the regime where cost matters. The distinction motivates a budget-indexed epistemic relation.

\begin{definition}[Resource-indexed knowledge]
For a resource scale $r\in\mathbb N$, write
\begin{equation}
K_i^{\le r}\varphi
\end{equation}
for the claim that agent $i$ can acquire an admissible epistemic warrant for $\varphi$ using at most $r$ units of the frozen resource measure.
\end{definition}

The word ``warrant'' is intentionally architecture-dependent. In a proof lane it may be a verifier-accepted proof certificate; in an experimental routing layer it may instead refer to a formally specified evidence state. The verifier boundary must be stated separately. In particular, a high-confidence forecast is not automatically $K_i^{\le r}\varphi$.

For a finite benchmark family $F$, one may measure bounded closure by
\begin{equation}
\Omega_i(r;F)=
\frac{
\left|\{\varphi\in F:T\vdash\varphi\text{ and }K_i^{\le r}\varphi\}\right|
}{
\left|\{\varphi\in F:T\vdash\varphi\}\right|
},
\end{equation}
when the denominator is nonzero. Ordinary logical omniscience corresponds to the ideal in which deductive consequence is available without a computationally meaningful resource barrier. DATA MIND instead studies the growth of accessible, certifiable consequence as resources increase.

\section{The nonstandard horizon}

Assume a nonstandard extension containing the hypernaturals ${}^*\mathbb N$. An element
\begin{equation}
H\in{}^*\mathbb N\setminus\mathbb N
\end{equation}
is called unlimited when it exceeds every standard natural number. There is no largest hypernatural: $H+1$ is larger than $H$. Consequently the construction below fixes an unlimited horizon rather than invoking a nonexistent ``largest infinity.''

Let
\begin{equation}
\varphi_0,\varphi_1,\varphi_2,\ldots
\end{equation}
be a standard enumeration of the formulas or sentences relevant to the chosen formal environment. Its nonstandard extension yields a hyperfinite initial segment
\begin{equation}
F_H=\{\varphi_n:n\le H\}.
\label{eq:FH}
\end{equation}
Internally, $F_H$ behaves like a finite family under transferred finite statements. Externally, it contains every standard-indexed $\varphi_n$ because $n<H$ for every standard $n$.

\subsection{What transfer does and does not provide}

The intended use of transfer is conservative: begin with a genuinely finitary, internal definition, and transfer that definition to hyperfinite indices. This permits finite combinatorial principles to be used on internally hyperfinite objects. It does \emph{not} justify replacing arbitrary external countable conjunctions by internal hyperfinite conjunctions without further work.

This distinction is essential. The expression
\begin{equation}
\bigwedge_{n\le H}{}^*\varphi_n
\end{equation}
is an internal hyperfinite conjunction, whereas
\begin{equation}
\bigwedge_{n<\omega}\varphi_n
\end{equation}
is a genuinely infinitary expression in an infinitary logic. The two constructions may be related, but their identification is not automatic.

\section{Hyperfinite omniscience}

We now state the object that motivates DATA MIND 3.2.

\begin{definition}[Hyperfinite omniscience at horizon $H$]
Let $T$ be the frozen theory and let $H$ be unlimited. Define
\begin{equation}
\Omn_i(H):=
\bigwedge_{\substack{n\le H\\T\vdash\varphi_n}}
{}^*K_i^{\le H}\varphi_n.
\label{eq:omn}
\end{equation}
\end{definition}

The statement says that every theorem encountered in the internal hyperfinite horizon is available to agent $i$ within the same unlimited resource horizon. It is a hyperfinite analogue of deductive closure, not a claim that a real machine has been built that literally knows all mathematics.

\subsection{Intersection semantics}

When propositions are represented by their truth sets over a space of worlds $W$, conjunction corresponds to set intersection:
\begin{equation}
[\![\varphi\wedge\psi]\!]=[\![\varphi]\!]\cap[\![\psi]\!].
\end{equation}
The hyperfinite analogue of \eqref{eq:omn} therefore has the truth-set representation
\begin{equation}
[\![\Omn_i(H)]\!]
=
\bigcap_{\substack{n\le H\\T\vdash\varphi_n}}
[\![{}^*K_i^{\le H}\varphi_n]\!].
\label{eq:intersection}
\end{equation}
This is the appropriate place for the intersection proposed in the motivating discussion: logical ``and'' becomes meet/intersection. By contrast, a supremum or maximum belongs naturally to the aggregation of \emph{costs}, not to Boolean truth values themselves.

\section{The omniscience horizon cost $\Ocal_i(H)$}

Suppose the least standard resource cost of acquiring a warranted settlement of $\varphi$ is
\begin{equation}
c_i(\varphi)=\min\{r\in\mathbb N:K_i^{\le r}\varphi\},
\end{equation}
whenever such a minimum is defined. Transfer produces the corresponding internal cost relation.

\begin{definition}[Hyperfinite epistemic closure cost]
For a fixed unlimited $H$, define
\begin{equation}
\boxed{
\Ocal_i(H)=
\max_{\substack{n\le H\\T\vdash\varphi_n}}
{}^*c_i(\varphi_n)
}
\label{eq:OH}
\end{equation}
when the internal cost is defined for the required hyperfinite family.
\end{definition}

The type of $\Ocal_i(H)$ is important:
\begin{equation}
\Ocal_i(H)\in{}^*\mathbb N.
\end{equation}
It is a hypernatural resource value, not a set of formulas and not a truth value. Since the indexed family is internally hyperfinite, the transferred finite maximum principle permits ``max'' rather than merely ``sup.''

A compact way to express hyperfinite closure within the chosen horizon is
\begin{equation}
\Ocal_i(H)\le H.
\label{eq:Homn}
\end{equation}
When \eqref{eq:Homn} holds, every required item in the horizon has cost no greater than the horizon budget, under the frozen cost definition.

\subsection{Normalized costs and shadows}

The hypernatural $\Ocal_i(H)$ may itself be unlimited and therefore need not have a finite real standard part. A normalized ratio may be limited:
\begin{equation}
R_i(H)=\frac{\Ocal_i(H)}{H}.
\end{equation}
If $R_i(H)$ is limited, define its standard-part shadow
\begin{equation}
\rho_i(H)=\st\!\left(\frac{\Ocal_i(H)}{H}\right)\in\mathbb R.
\label{eq:rho}
\end{equation}
This quantity is a standard real summary of normalized hyperfinite closure cost. It should not be confused with the ``standard trace'' of a set.

\section{Standard traces and standard-part shadows}

The word ``shadow'' is most precise for a limited hyperreal or related nearstandard object. For a hyperfinite set such as $F_H$, the safer set-theoretic language is its \emph{standard trace}: the members whose indices or elements are standard.

For example,
\begin{equation}
F_H^{\mathrm{st}}
=
\{\varphi_n:n\in\mathbb N\}
\end{equation}
in the sense that every standard-indexed formula is present below any unlimited $H$.

Now define the internal successful-theorem index set
\begin{equation}
E_i(H)=
\{n\le H:T\vdash\varphi_n\text{ and }{}^*K_i^{\le H}\varphi_n\}.
\end{equation}
Its standard trace records which standard-indexed theorems lie inside the hyperfinite epistemic state. If one instead converts the hyperfinite set to a bounded hyperreal density,
\begin{equation}
\Omega_i(H)=
\frac{|E_i(H)|}{|\{n\le H:T\vdash\varphi_n\}|},
\end{equation}
then, when the denominator is nonzero,
\begin{equation}
\omega_i(H)=\st(\Omega_i(H))\in[0,1]
\label{eq:shadowcoverage}
\end{equation}
is a genuine standard-part shadow of hyperfinite coverage density.

Thus two different summaries coexist:
\begin{align}
\omega_i(H)&=\text{shadow of hyperfinite coverage density},\\
\rho_i(H)&=\text{shadow of normalized worst-case closure cost}.
\end{align}
They answer different questions. High coverage does not imply low worst-case cost, and a low worst-case ratio says nothing by itself about the quality of the settlement role chosen.

\section{From theorem knowledge to settlement knowledge}

DATA MIND is not intended to be only a theorem prover. In the AMLD view, a target can be approached through distinct settlement roles. Let
\begin{equation}
\mathcal S=\{P,R,I,C,U\},
\end{equation}
where $P$ denotes proof-directed settlement, $R$ refutation-directed settlement, $I$ independence-directed settlement, $C$ contradiction or integrity-directed settlement, and $U$ unresolved or withheld status.

The distinction matters because a system can waste enormous resources performing excellent proof search on a problem whose correct diagnostic role is not proof. The epistemic question is therefore not only
\begin{quote}
Does the system possess enough search power to produce a certificate?
\end{quote}
but also
\begin{quote}
Does the system recognize what kind of certificate should be sought, and allocate attention accordingly?
\end{quote}

This is the point at which an oracle becomes scientifically useful as a diagnostic idealization.

\section{The hyperfinite settlement oracle}

\begin{definition}[Hyperfinite settlement oracle]
For the hyperfinite horizon $F_H$, define an idealized internal map
\begin{equation}
\boxed{
\mathscr O_H:F_H\longrightarrow\{P,R,I,C,U\}.
}
\label{eq:oracle}
\end{equation}
For each target in its declared domain, $\mathscr O_H$ records the gold settlement role under the frozen semantics and provenance rules.
\end{definition}

The phrase ``oracle'' is intentionally conditional. In mathematics one may study a hypothetical oracle whose correctness is included in its specification. That is not the same as claiming that an implemented classifier is actually omniscient. DATA MIND 3.2 therefore separates the ideal mathematical oracle from finite executable records whose ground truth must have independent provenance.

\subsection{Four possible oracle strengths}

It is useful to distinguish several oracle interfaces.

\begin{enumerate}[label=\textbf{O\arabic*.},leftmargin=2.3em]
\item \textbf{Role oracle.} Returns only the correct member of $\{P,R,I,C,U\}$.
\item \textbf{Resource oracle.} Returns an ideal allocation or recipient of marginal search resources.
\item \textbf{Strategy oracle.} Returns an ideal next search direction or action class.
\item \textbf{Certificate oracle.} Returns a purported settlement certificate $\pi$.
\end{enumerate}

These strengths should not be conflated. A role oracle removes role-identification uncertainty but still leaves the actual settlement problem intact. A certificate oracle is much stronger; even then, DATA MIND's verifier boundary should remain in force.

\section{How the oracle enters AMLD}

DATA MIND 3.2 uses the oracle as an epistemic instrument, not as a sovereign mathematical authority. The permitted flow is
\begin{equation}
\boxed{
\mathscr O
\longrightarrow
\text{hidden evaluation or awareness/routing cue}
\longrightarrow
P/R/I/C\text{ search}
\longrightarrow
\V
\longrightarrow
\BANK.
}
\label{eq:flow}
\end{equation}
The forbidden shortcuts are
\begin{equation}
\mathscr O\not\longrightarrow\BANK,
\qquad
\mathscr O\not\longrightarrow\text{verifier acceptance}.
\label{eq:forbidden}
\end{equation}

\begin{principle}[Verifier sovereignty]
An oracle output may change what a search agent attends to or attempts, but it cannot by itself promote a mathematical claim into certified shared knowledge.
\end{principle}

In compact form:
\begin{center}
\emph{The oracle may direct attention; it may not manufacture certified knowledge.}
\end{center}

That principle preserves the central epistemic asymmetry of DATA MIND. Search components may be probabilistic, reflective, creative, or counterfactual. Certification remains exact under the frozen verifier policy.

\section{The three DATA MIND 3.2 oracle modes}

The current 3.2 architecture defines three distinct scientific modes. Their separation is necessary for interpretable experiments.

\subsection{Mode A: hidden ground truth}

In the default mode the oracle is invisible to the AMLD society during search. The system receives only the formal problem. After search, an evaluator may compare the reported settlement role with the frozen oracle record.

If $s_n$ is AMLD's reported role and $o_n=\mathscr O_H(\varphi_n)$ is the gold role, define
\begin{equation}
E_H=\{n\le H:s_n=o_n\}.
\end{equation}
A hyperfinite role-accuracy quantity is then
\begin{equation}
A_H=\frac{|E_H|}{H+1},
\end{equation}
and, when interpreted in a suitable nonstandard model,
\begin{equation}
\st(A_H)
\end{equation}
would be a standard shadow of role accuracy. In ordinary executable experiments, DATA MIND uses only the finite analogue of this calculation.

\subsection{Mode B: Professor role hint}

In the second mode the role oracle is opened only far enough to change awareness. If the gold role is $I$, for example, the architecture may make independence-directed work salient to the Professor-facing independence agent $I1$. The protected partner $I2$ remains unhinted.

This mode does not assert the target theorem, its negation, or an independence certificate. It merely supplies the role label. The Professor remains an adviser, not a source of theoremhood.

The counterfactual question is:
\begin{quote}
How much settlement cost is attributable to having to recognize the correct epistemic role before doing the corresponding search?
\end{quote}

\subsection{Mode C: direct role hint}

The stronger counterfactual gives the correct role to both members of the matching couple. This deliberately removes more of the routing problem and therefore must be labeled oracle-assisted performance rather than ordinary DATA MIND performance.

The distinction between Modes B and C also measures the value of the protected independent lane. If Professor-only guidance helps while direct guidance harms, the result would be evidence that diversity and anti-herding protection matter even under perfect role information.

\section{Professor, awareness, and oracle information}

The Professor in DATA MIND is best interpreted as an awareness and attention mechanism. It need not be logically omniscient. The oracle can, in an experimental arm, provide a perfect role label to the Professor-facing bridge, but this does not convert the Professor into an infallible mathematician.

Let $A_i\varphi$ represent that formula or direction $\varphi$ is within agent $i$'s operative awareness. Then an oracle role hint is naturally modeled as an intervention on awareness or routing:
\begin{equation}
\mathscr O_H(\tau)=I
\quad\leadsto\quad
A_{I1}(\text{independence-directed search on }\tau),
\end{equation}
not as the epistemic assertion
\begin{equation}
K_{I1}(\text{``$\tau$ is independent''}).
\end{equation}
The latter would require a justified epistemic warrant under the declared semantics. This separation avoids confusing salience with knowledge.

\section{BANK, FUTUREBANK, and the oracle boundary}

The AMLD memory architecture gives the oracle a natural place without weakening trust.

\begin{description}[leftmargin=2.1cm,style=nextline]
\item[$\BANK$] Stores material admitted under the frozen verification policy. It represents shared certified mathematical capital.
\item[$\FBANK$] Stores speculative, counterfactual, or prospective material. It may contain attractive possibilities that are not certified.
\item[Oracle] Provides hidden gold information or experimental guidance. Oracle outputs do not become BANK deposits merely because the oracle is idealized.
\end{description}

A certificate-producing oracle, if studied, should therefore satisfy the same pattern as any other proposer:
\begin{equation}
\mathscr O^{\mathrm{cert}}(T,\varphi)=\pi,
\qquad
\V(T,\varphi,\pi)=1
\end{equation}
before $\pi$ contributes certified knowledge. This preserves a clean distinction between the \emph{source} of a candidate and the \emph{authority} that accepts it.

\section{Role-specific omniscience profiles}

One attraction of AMLD is that logical omniscience need not collapse to a single scalar. The system may have very different abilities in the four settlement directions. For a frozen role-labeled horizon, define role-restricted closure costs
\begin{equation}
\Ocal_P(H),\quad
\Ocal_R(H),\quad
\Ocal_I(H),\quad
\Ocal_C(H),
\end{equation}
where each maximum is taken only over targets whose correct settlement role is the corresponding component.

Likewise define coverage shadows, when meaningful,
\begin{equation}
\omega_P(H),\quad
\omega_R(H),\quad
\omega_I(H),\quad
\omega_C(H).
\end{equation}
An architecture might be near closure on ordinary proof problems while weak at detecting contradictions or constructing independence certificates. This would expose the ``not-ATP'' part of conjecture settlement mathematically rather than rhetorically.

A role-confusion matrix adds directional information. For true role $a$ and reported role $b$, define the hyperfinite count
\begin{equation}
M_{ab}(H)=
|\{n\le H:\mathscr O_H(\varphi_n)=a,\;\AMLD(\varphi_n)=b\}|.
\end{equation}
Finite experiments can compute the ordinary analogue exactly.

\section{Oracle gaps as diagnostic upper bounds}

Let $c_{\mathrm{normal}}(\varphi)$ denote the observed resource cost without role information and let $c_{\mathrm{role}}(\varphi)$ denote the cost when correct role information is supplied under a frozen counterfactual interface. Aggregate over a finite or hyperfinite family in a common way.

At a finite executable horizon $N$, a simple worst-case diagnostic is
\begin{equation}
\Delta_N=
\Ocal_{\mathrm{normal}}(N)-
\Ocal_{\mathrm{role}}(N),
\label{eq:deltaN}
\end{equation}
when both quantities are defined.

A large positive $\Delta_N$ suggests that role identification or routing consumes a substantial part of settlement cost. A small $\Delta_N$ suggests that the dominant difficulty remains inside the selected role: proof search, countermodel construction, independence certification, contradiction localization, presentation choice, or some other internal bottleneck.

The oracle gap is therefore a \emph{counterfactual diagnostic upper bound}, not evidence that a realizable perfect oracle exists. It asks how much performance could improve if one particular uncertainty were removed while the rest of the architecture remained frozen.

\section{Finite executable proxies in DATA MIND 3.2}

A conventional computer cannot literally instantiate an unlimited hypernatural. DATA MIND 3.2 therefore implements a finite proxy rather than pretending otherwise.

For a finite frozen horizon
\begin{equation}
F_N=\{\varphi_n:n\le N\},
\end{equation}
the implementation may compute:
\begin{itemize}
\item finite settlement coverage;
\item the finite maximum $\Ocal_i(N)$ when every required least observed cost is available;
\item the normalized ratio $\Ocal_i(N)/N$ when defined;
\item role confusion matrices and false-positive counts;
\item hidden-ground-truth, Professor-role-hint, and direct-role-hint comparisons;
\item finite oracle gaps such as \eqref{eq:deltaN}.
\end{itemize}

These are empirical approximations to the mathematical structure. They are not reported as literal values of $\st(\Omega_i(H))$ or $\st(\Ocal_i(H)/H)$.

The distinction is more than terminological. It prevents an implementation convenience from being mistaken for a nonstandard theorem. The mathematical layer provides a conceptual limiting geometry; the executable layer supplies finite evidence.

\section{A motivating first experiment: the PRCOM epistemic role quartet}

The first proposed DATA MIND 3.2 experiment is naturally built around the Metamath theorem \texttt{prcom}, whose statement is the unordered-pair commutativity identity
\begin{equation}
\{A,B\}=\{B,A\}.
\end{equation}
PRCOM has already been a revealing proof-search target because the system can encounter severe branching and target-drift behavior. A single ordinary PRCOM run, however, exercises only the $P$ role. The oracle layer is better tested on a matched family whose epistemic status varies while the mathematical nucleus remains recognizable.

The proposed quartet is:
\begin{center}
\begin{tabular}{p{0.18\textwidth}p{0.58\textwidth}c}
\toprule
Case & Construction & Gold role\\
\midrule
PRCOM-P & Ordinary \texttt{prcom} in the normal consistent theory & $P$\\
PRCOM-R & A target representing the negation of PRCOM in the same consistent environment & $R$\\
PRCOM-C & A supplied context containing PRCOM and its negation, with an otherwise ordinary target & $C$\\
PRCOM-I & A deliberately constructed environment in which a related target is genuinely independent, supported by an appropriate two-world/model certificate & $I$\\
\bottomrule
\end{tabular}
\end{center}

The independence case requires special care. Removing an axiom and then failing to prove the target does not establish independence. A defensible $I$ case should provide certified models or equivalent witnesses showing compatibility with both the target and its negation under the frozen base theory.

A three-arm comparison can then hold problems, seeds, verifier, budgets, and stopping rules fixed:
\begin{enumerate}
\item hidden oracle only;
\item Professor-facing role hint;
\item direct role hint to the matching couple.
\end{enumerate}

Primary outcomes would include verifier-valid settlement, false positives, role confusion, resources by role, and oracle-gap quantities. This note records the design motivation only. It does \textbf{not} launch or preregister DATA MIND 3.2 Experiment 1.

\section{What would count as evidence?}

The oracle experiment is useful because several qualitatively different outcomes are possible.

\begin{enumerate}[label=\arabic*.]
\item \textbf{Large improvement from role hints.} Role recognition and resource routing are major bottlenecks. Work should focus on settlement classification, awareness, Professor calibration, couple disagreement, and role-switching signals.
\item \textbf{Little improvement from role hints.} The dominant bottleneck lies inside the selected role. For PRCOM-P, this would point back toward branching control, target relevance, presentation trading, Quotient Hunter, macros, depth, or other proof-search mechanisms.
\item \textbf{Professor-only beats direct hint.} The protected independent partner lane adds value even when the role signal is idealized, supporting the anti-herding design.
\item \textbf{Direct hint beats Professor-only strongly.} The Professor-to-agent communication channel, not role information itself, may be the limiting layer.
\item \textbf{Oracle hint harms performance.} Even correct high-level information can be operationally destructive if the mechanism translating it into search control oversteers the system.
\end{enumerate}

The last possibility is especially important in light of earlier DATA MIND Professor experiments: verifier safety does not imply strategic safety. A bad control response can waste resources without ever causing the verifier to accept a false certificate.

\section{Open mathematical questions}

The hyperfinite formulation creates several research directions that should be kept separate from the executable engineering claims.

\subsection{Dependence on enumeration}

Both hyperfinite density and horizon costs may depend on the chosen enumeration of formulas. A robust theory should state invariance conditions or compare equivalence classes of admissible enumerations rather than treating one coding as canonical without justification.

\subsection{Dependence on presentation}

Proof cost is presentation-sensitive. Certified presentation trading may change $c_i(\varphi)$ and therefore $\Ocal_i(H)$ without changing provability. This gives a natural meeting point between the hyperfinite oracle framework and the trading-theorem program: semantic content may be invariant while the finite or hyperfinite cost geometry changes.

\subsection{Hyperfinite versus infinitary closure}

An internal hyperfinite conjunction is not automatically a formula of an external infinitary logic, and a theorem about one should not be silently reported as a theorem about the other. Establishing precise bridge results would be mathematically valuable.

\subsection{Standard extraction}

If a standard target is settled by an internal hyperfinite derivation of unlimited length, additional work is needed before one may infer the existence of an appropriate standard finite certificate. Transfer alone does not license that collapse. Standard extraction, underspill, compactness, or a system-specific argument may be required, depending on the formalization.

\subsection{Oracle provenance}

A finite benchmark oracle is only as trustworthy as its labels. $P$ and $R$ cases may often be grounded by certificates. $I$ and $C$ cases demand equally explicit semantics and witnesses. ``Gold label'' must never become a synonym for ``label we decided to trust.''

\section{Relation to epistemic axioms}

The oracle construction does not require the implemented AMLD agents themselves to satisfy the full S5 package. Several modal properties should be distinguished.

\begin{align*}
\mathbf K:&\quad K_i(\varphi\to\psi)\to(K_i\varphi\to K_i\psi),\\
\mathbf T:&\quad K_i\varphi\to\varphi,\\
\mathbf 4:&\quad K_i\varphi\to K_iK_i\varphi,\\
\mathbf 5:&\quad \neg K_i\varphi\to K_i\neg K_i\varphi.
\end{align*}

For DATA MIND, factivity is most naturally attached to verifier-certified knowledge, while positive and negative introspection concern self-awareness. Full logical omniscience arises from closure assumptions such as normality and necessitation; it should not be inferred merely from strong self-awareness. A practical architecture may therefore be factive at the certification boundary, introspectively bounded, and computationally far from deductive closure.

The hyperfinite oracle provides an ideal comparison object without forcing AMLD itself to be an ideal epistemic agent.

\section{Architecture summary}

The conceptual division of labor in DATA MIND 3.2 can be written compactly:
\begin{center}
\begin{tabular}{p{0.22\textwidth}p{0.66\textwidth}}
\toprule
Component & Epistemic role\\
\midrule
Hyperfinite oracle & Idealized gold information for mathematical analysis and controlled counterfactuals\\
Professor & Changes salience, attention, and advisory routing; does not certify theoremhood\\
$P/R/I/C$ agents & Attempt the corresponding forms of settlement\\
Verifier $\V$ & Determines whether a candidate certificate satisfies the frozen acceptance rule\\
$\BANK$ & Stores shared verifier-admitted mathematical capital\\
$\FBANK$ & Stores speculative and counterfactual possibilities without promoting them to truth\\
Sentinel & May veto resource-pathological actions; does not determine mathematical truth\\
\bottomrule
\end{tabular}
\end{center}

The architecture therefore allows a deliberately powerful epistemic idealization while keeping the trusted computing boundary narrow.

\section{Conclusion}

Logical omniscience is usually presented as an unrealistic consequence of ideal epistemic closure. DATA MIND 3.2 turns that limitation into a measurable research object. The operator $K_i^{\le r}$ makes resource cost explicit. An unlimited hypernatural horizon $H$ yields a conceptual hyperfinite closure problem, and $\Ocal_i(H)$ records a hypernatural worst-case closure cost. Standard traces recover standard-indexed behavior, while normalized limited hyperreals admit standard-part shadows.

The hyperfinite oracle then moves the construction from pure epistemic semantics into AMLD. By selectively withholding or revealing only the correct settlement role, one can measure the value of perfect role information without handing the system a certificate. This creates a controlled diagnostic of the ``not-ATP'' part of settlement: recognition, routing, awareness, and coordination among proof, refutation, independence, and contradiction agents.

The decisive trust principle remains unchanged. Oracle information may be idealized, Professor advice may be intelligent, and search may be creative, but none of them are permitted to manufacture certified mathematics. The verifier remains sovereign, and BANK remains the shared record of what has actually crossed that boundary.

\appendix
\section{Notation}

\begin{center}
\begin{tabular}{p{0.24\textwidth}p{0.66\textwidth}}
\toprule
Symbol & Meaning\\
\midrule
$K_i\varphi$ & Agent $i$ knows $\varphi$ in the selected epistemic semantics\\
$K_i^{\le r}\varphi$ & Agent $i$ can acquire the declared warrant for $\varphi$ within budget $r$\\
$H$ & Fixed unlimited hypernatural used in the mathematical specification\\
$F_H$ & Internal hyperfinite formula horizon $\{\varphi_n:n\le H\}$\\
$\Omn_i(H)$ & Hyperfinite conjunction of required epistemic closures through $H$\\
$c_i(\varphi)$ & Least resource cost, when defined\\
$\Ocal_i(H)$ & Hyperfinite maximum closure cost through horizon $H$\\
$\rho_i(H)$ & Standard part of $\Ocal_i(H)/H$ when limited\\
$\omega_i(H)$ & Standard part of hyperfinite coverage density when defined\\
$\mathscr O_H$ & Hyperfinite settlement-role oracle\\
$P,R,I,C,U$ & Proof, refutation, independence, contradiction, unresolved/withheld\\
$\V$ & External/frozen verifier boundary\\
$\BANK$ & Shared verifier-admitted mathematical capital\\
$\FBANK$ & Speculative transactional future state\\
\bottomrule
\end{tabular}
\end{center}

\section{Implementation-status statement}

As of September 2026, DATA MIND 3.2 contains a reversible finite-horizon oracle research layer on a separate development branch. Its executable classes are finite proxies of the nonstandard specification. The implementation distinguishes hidden evaluation, Professor-facing role hints, and direct role hints, and it marks oracle-derived cues as non-certified. The layer has interface tests but, at the time of this note, DATA MIND 3.2 Experiment 1 has not been launched.

\begin{thebibliography}{9}

\bibitem{Hintikka1962}
J. Hintikka,
\emph{Knowledge and Belief: An Introduction to the Logic of the Two Notions},
Cornell University Press, 1962.

\bibitem{Fagin1995}
R. Fagin, J. Y. Halpern, Y. Moses, and M. Y. Vardi,
\emph{Reasoning About Knowledge},
MIT Press, 1995.

\bibitem{Robinson1966}
A. Robinson,
\emph{Non-standard Analysis},
North-Holland, 1966.

\bibitem{Goldblatt1998}
R. Goldblatt,
\emph{Lectures on the Hyperreals: An Introduction to Nonstandard Analysis},
Springer, 1998.

\bibitem{ChangKeisler1990}
C. C. Chang and H. J. Keisler,
\emph{Model Theory}, 3rd ed.,
North-Holland, 1990.

\bibitem{BarwiseFeferman1985}
J. Barwise and S. Feferman, editors,
\emph{Model-Theoretic Logics},
Springer, 1985.

\end{thebibliography}

\end{document}
